\section{Proofs of the root bounds and additional boundary cases}
\label{sec:gamma-annulus-proof}

\subsection{Bounds for a Gamma polynomial}

For \(\gamma>0\) put
\[
 G_{m,\gamma}(z)=\sum_{k=0}^m\binom{m}{k}(\gamma)_kz^{m-k},
 \qquad h_0=m-1 .
\]
We first show that every zero \(z\) of \(G_{m,\gamma}\) satisfies
\[
 |z|\le\gamma+h_0,\qquad
 \Re z\le-\gamma+h_0.
\]
Let \(T=\gamma+m\). The polynomial is \(\det(zI+J)\), where
\(J=D+iK\) is complex symmetric, with real diagonal
\(D_{jj}=T+1-2j\) for \(1\le j\le m\) and real symmetric
off-diagonal entries \(K_{j,j+1}=\sqrt{j(T-j)}\).
To check the representation, the characteristic polynomial of the
leading \(j\times j\) block is
\[
 \sum_{k=0}^j(-1)^{j-k}\frac{j!}{k!}\binom{T-1-k}{j-k}z^k .
\]
This expression has the initial terms \(1,z-(T-1)\) and
satisfies the determinant recurrence, whose two coefficients are
\(z-(T+1-2j)\) and \((j-1)(T+1-j)\). At \(j=m\), replacing \(z\)
by \(-z\) and multiplying by \((-1)^m\) gives \(G_{m,\gamma}\).

The lower bound on the eigenvalue real parts of \(J\) follows
from its real diagonal \(D\). To bound its numerical radius,
take \(0<\theta<\pi\) and set \(q=\tan(\theta/2)\).
For the Hermitian part \(H_\theta=(\cos\theta)D+(\sin\theta)K\) of
\(e^{-i\theta}J\), the positive vector with entries
\[
 v_j=\sqrt{\binom{T-1}{j-1}}\;q^{j-1}
\]
satisfies \((H_\theta v)_j=(T-1)v_j\) for \(j<m\), and
\((H_\theta v)_m=[T-1-\gamma(1-\cos\theta)]v_m\).
These identities follow by substituting the adjacent ratios
in the off-diagonal terms. The top eigenvector of \(H_\theta\)
can be taken positive by maximizing its Rayleigh quotient
with nonnegative coordinates and using the connected off-diagonal
chain. Its scalar product with the displayed inequality gives
\(\lambda_{\max}(H_\theta)<T-1\). An alternating diagonal
conjugates \(K\) to \(-K\), and continuity handles the endpoint
angles. Hence \(|z|\le T-1\) for each eigenvalue \(z\) of
\(-J\). For \(m\ge2\), equality could only occur in an endpoint
direction. An eigenvector would then be supported in a single
extreme diagonal coordinate, which the adjacent entry prevents.
Thus the modulus bound is strict for \(m\ge2\).

Coefficient comparison also gives the equation
\[
 zG''+(1-\gamma-m-z)G'+mG=0 .
\]
The constant term is positive. At a repeated nonzero zero the
term \(zG''\) would have strictly smaller order than every
other term, so the zeros are simple. Let \(z_i\) minimize
the modulus and write \(r=|z_i|\). The real part of the
differential equation at \(z_i\), using
\(G''(z_i)/G'(z_i)=2\sum_{j\ne i}(z_i-z_j)^{-1}\), gives
\begin{equation}\label{eq:gamma-radius-sum}
 r-\gamma\ \ge\ -\Re z_i-\gamma
  =\sum_{j\ne i}\frac{|z_j|^2-r^2}{|z_j-z_i|^2}\ \ge\ 0 .
\end{equation}
Equality \(r=\gamma\) forces \(z_i=-\gamma\) and all moduli
\(\gamma\). Since their sum is \(-m\gamma\), all zeros would
then equal \(-\gamma\). Thus \(r>\gamma\) for \(m\ge2\).

When \(m\ge2\) and \(\gamma\ge15h_0\), we show the stronger
\(r>\gamma+h_0/6\). Vieta's formulas give
\(\sum z_j=-m\gamma\) and
\(\sum z_j^2=m\gamma^2-mh_0\gamma\). In particular,
\[
 \sum_j |z_j|^2\ge m\gamma^2+mh_0\gamma .
\]
Set \(u_j=z_j+\gamma\) and \(s_0=h_0/\gamma\).
The two bounds on \(z_j\) above give \(|\Re u_j|\le h_0\).
Since \(\sum u_j=0\) and \(\sum u_j^2=-mh_0\gamma\), it
follows that \(\sum |u_j|^2\le mh_0\gamma(1+2s_0)\).
Suppose \(r\le\gamma+h_0/6\), and set
\(\eta=(r^2-\gamma^2)/(\gamma h_0)\). Then
\(0\le\eta\le181/540\) and
\(|u_i|^2\le r^2-\gamma^2\), since \(\Re z_i\le-\gamma\).
Write \(q_j=|z_j|^2-r^2\). We have
\[
 \sum_j q_j\ge m\gamma h_0(1-\eta),\qquad
 0\le q_j\le\gamma h_0(2+s_0-\eta),\qquad
 \sum_j |z_j-z_i|^2\le m\gamma h_0(1+2s_0+\eta).
\]
The last inequality follows from \(\sum u_j=0\).
Weighted Cauchy--Schwarz in \eqref{eq:gamma-radius-sum} now gives
\[
 r-\gamma\ge
 m\frac{(1-\eta)^2}{(2+s_0-\eta)(1+2s_0+\eta)}
 \ge m\,\frac{128881}{741455}>\frac{h_0}{6}.
\]
The rational expression decreases in \(s_0\) and \(\eta\)
on the stated ranges, so its value at \(s_0=1/15,\eta=181/540\)
gives the middle bound. This contradiction proves the stronger
inner radius.

\subsection{Products of Gamma annuli}

For monic polynomials of degree \(m\) in the binomial basis, write
\[
 \left(\sum_{k=0}^m\binom{m}{k}a_k z^{m-k}\right)\diamond
 \left(\sum_{k=0}^m\binom{m}{k}b_k z^{m-k}\right)
 =\sum_{k=0}^m\binom{m}{k}a_kb_k z^{m-k}.
\]
The upper and lower root radii multiply under this operation.
For the upper bound, if the first polynomial is
\(\prod_{j=1}^m(z+\delta_j)\) and \(Q^{\mathrm{pol}}\) is the
polarization of the second, the composed value at \(z\) equals
\(\prod_j\delta_j\, Q^{\mathrm{pol}}(z/\delta_1,\ldots,z/\delta_m)\).
If \(|z|\) exceeds the product of the two upper radii, all the
arguments lie in the open exterior of the disk for the second
polynomial. The Grace--Walsh--Szeg\H{o} theorem
\cite[Theorem~4.1]{wagner-stability} gives nonvanishing there.
The lower bound follows by reversal, which commutes with
\(\diamond\). With strict radii one can first tighten the
bounds on the finitely many roots of each factor.

Use \(\gamma=am+1\) and the identity
\[
 (\gamma)_{dk}=d^{dk}\prod_{j=0}^{d-1}((\gamma+j)/d)_k .
\]
It expresses \(R_m\), after scaling the variable by \(d^d\),
as the composition of the \(d\) Gamma polynomials of shapes
\((\gamma+j)/d\). Consequently
\[
 (\gamma)_d\le |x|\le(\gamma+d(m-1))_d
 \qquad(R_m(x)=0),
\]
strictly for \(m\ge2\). This gives the first assertion of
Lemma~\ref{lem:gamma-annulus}. At the two orders
\(\gamma=am+1\), the inner radius for \(m=N\) is at least
the outer radius for \(m=N-1\) when \(a\ge d(N-2)\).
The strict inner bound gives separation even at equality and
at \(N=2\). If \(N\ge25\) and \(6a\ge(5N-11)d\), then
\(a\ge15d\). Each of the \(d\) unit-step shapes for
\(R_N\) is at least \(15(N-1)\). The stronger bound gives
the lower radius
\((aN+1+d(N-1)/6)_d\), and
\[
 aN+1+d(N-1)/6\ge a(N-1)+1+d(N-2).
\]
This proves the remaining separation in
Lemma~\ref{lem:gamma-annulus}.

\subsection{Small positive lower exponents and adjacent exponents}

\begin{proposition}
\label{prop:additional-boundary}
For \(N\ge25\), the pair \(P_{N-1},P_N\) is coprime on
each of the following supports:
\[
 1\le a\le4,\quad 1\le d<3a,
 \qquad\text{or}\qquad
 a\ge1,\quad d=1.
\]
\end{proposition}

\begin{proof}
Suppose first \(N\ge3150\). We may assume
\(6a<(5N-11)d\), by Lemma~\ref{lem:gamma-annulus}. In both
domains \(d<3a\), and Lemma~\ref{lem:degree-cap} applies
because \(aN\ge1375\). Thus the primes of the common-factor
norm in Lemma~\ref{lem:resultant} are supported up to \(C_N\)
as in Lemma~\ref{lem:prime-support}.

If \(1\le a\le4\) and \(d<3a\), then \(b\le15\),
\(a+2d-1\le25\), and \(s\le121/15\). Use the prime threshold
\(X=\sqrt{15N}\), which is at least \(C_N\) and exceeds \(150\).
By Lemmas~\ref{lem:prime-estimates} and~\ref{lem:small-prime-sums},
the right side of \eqref{eq:pade-norm} is less than
\[
 \tfrac12(\log N+\log 15)-.45+2.48\sqrt{15/N}.
\]
Here \(\log(M+1)/\log X<2.001\), and
\(2.001(1+1.2762/\log X)<2.48\).
The loss bound \(g\le7s/32\) in \eqref{eq:pade-loss} gives
\(\log\Lambda>\log N-2.776\).
The difference exceeds \(0.17\) at \(N=3150\) and increases
with \(N\).

For \(d=1\) and \(a\ge2\) take \(X=\sqrt{(a+1)N}=C_N\),
since \(a+1<(5/6)(N-1)\). The prime-sum and prime-count
bounds give instead
\[
 \tfrac12(\log N+\log(a+1))+.2+2.60\sqrt{(a+1)/N}.
\]
Use the \(d=1\) gain from Lemma~\ref{lem:resultant}.
The surplus is greater than
\[
 \tfrac12\log N-.701-\tfrac12\log(3/2)
       -\log(1+1/\sqrt2)-2.60\sqrt{5/6}>.21
 \qquad(N\ge3150).
\]
The case \(a=d=1\) is in the preceding domain.

For \(25\le N<3150\), the supplemental checker
\texttt{verify\_small\_support.cpp} treats the two stated
support families together. It retains the inequalities
\(a<d(N-2)\), \(6a<(5N-11)d\) and \(d^2<19(a+d)\), using
Lemmas~\ref{lem:gamma-annulus} and~\ref{lem:ratio-range},
and uses the gap caps in Lemmas~\ref{lem:gap-caps}
and~\ref{lem:local-gap}. It groups the first factorial block
test by \((N,a)\) and then applies the individual block and
Newton slope tests of Section~\ref{sec:arithmetic}.
The resulting counts are \(12\,520\) groups and \(18\,904\)
supports. Of these, \(4\,357\) are excluded by the degree-dependent
prime bounds and \(4\,853\) by comparing Newton polygons.
For the other \(9\,694\), the full monic reciprocals are coprime
modulo \(1\,000\,000\,007\). The checker verifies primality by
trial division and computes the modular gcd by exact Euclidean
division. A monic common factor over \(\mathbb Q\) would be
integral and would retain its degree modulo this prime.
Thus these tests establish coprimality throughout the finite range.
\end{proof}

